\chapter{存储 I/O 栈、可靠性与端到端 Trace}

\begin{keyidea}
文件读取只有在文件层、块层、设备队列、DMA、介质和完成路径全部闭合后才真正结束。本章用端到端 trace 收束整个存储大部分。
\end{keyidea}

\section{一、读一个字节与写一个字节：四条端到端 trace}

前面六节分别讲了每种器件的内部结构。本节用同一个动作——CPU 读一个字节、写一个字节——把信号从 CPU 寄存器出发，逐级追踪到物理单元，再追踪回来。每条 trace 标注每一步的耗时，让读者看到延迟到底花在哪里。

\topic{8.1\quad SRAM：读一个字节（$\sim$1--2 ns）}

场景：CPU 执行 \code{MOV AL, [RSP+8]}，地址命中 L1D cache。

\begin{longtable}{L{0.06\linewidth}L{0.30\linewidth}L{0.30\linewidth}L{0.24\linewidth}}
\toprule
步 & 硬件事件 & 信号路径 & 耗时 \\
\midrule
1 & CPU 发出 load 地址 & 地址总线 $\to$ L1D tag 阵列 & 0 ns（组合逻辑） \\\tablerowrule
2 & Tag 比较命中 & 32 个 tag 并行比较 $\to$ hit 信号 & $\sim$0.3 ns \\\tablerowrule
3 & 行地址译码 & 地址 [11:6] $\to$ 行译码器 $\to$ WL[42] 拉高 & $\sim$0.2 ns \\\tablerowrule
4 & 字线激活 & WL[42] 驱动该行所有 6T 单元的 M5/M6 导通 & $\sim$0.1 ns \\\tablerowrule
5 & 位线放电/保持 & 8 对 BL 各自：Q=0 的单元通过 M5 放电 $\sim$80 mV；Q=1 的保持 VDD & $\sim$0.3 ns \\\tablerowrule
6 & 灵敏放大器锁存 & 8 个 SA 检测 BL 与 $\overline{BL}$ 的 $\pm$80 mV 差值，正反馈放大到全摆幅 & $\sim$0.4 ns \\\tablerowrule
7 & 列多路选择 & 地址 [5:3] 选择 8:1 列 mux 中的一组 8 bit & $\sim$0.1 ns \\\tablerowrule
8 & 数据返回 CPU & 8 bit 数据总线 $\to$ 对齐 $\to$ AL 寄存器 & $\sim$0.1 ns \\
\midrule
\multicolumn{3}{L{0.66\linewidth}}{总延迟（L1D 命中）} & $\sim$1.5 ns \\
\bottomrule
\end{longtable}

\begin{pdfdiagram}[x=1cm,y=1cm]
  \node[hw state, minimum width=1.4cm] (cpu) at (0,0) {CPU};
  \node[hw compute, minimum width=1.2cm] (tag) at (2.0,0) {Tag};
  \node[hw control, minimum width=1.2cm] (dec) at (4.0,0) {行译码};
  \node[hw state, minimum width=1.2cm] (wl) at (6.0,0) {WL};
  \node[hw compute, minimum width=1.2cm] (cell) at (8.0,0) {6T$\times$8};
  \node[hw compute, minimum width=1.2cm] (sa) at (10.0,0) {SA};
  \node[hw compute, minimum width=1.2cm] (mux) at (12.0,0) {列 mux};
  \draw[hw bus] (cpu) -- node[above,font=\tiny]{addr} (tag);
  \draw[hw bus] (tag) -- node[above,font=\tiny]{hit} (dec);
  \draw[hw bus] (dec) -- (wl);
  \draw[hw bus] (wl) -- node[above,font=\tiny]{WL=1} (cell);
  \draw[hw bus] (cell) -- node[above,font=\tiny]{$\pm$80mV} (sa);
  \draw[hw bus] (sa) -- node[above,font=\tiny]{full swing} (mux);
  \draw[hw bus] (mux) -- ++(1.0,0) node[right,font=\tiny]{8 bit $\to$ AL};
  \node[hw label, text width=10cm] at (6.0,-1.0) {SRAM 读 1 字节：地址→tag 命中→行译码→WL 打开→8 对 BL 放电→SA 放大→列 mux→数据返回。全程 $\sim$1.5 ns。};
\end{pdfdiagram}
\diagramnote{SRAM 读路径。每一步都是组合逻辑或单级放大，没有等待状态，因此 L1 访问可以在一个时钟周期内完成。}

\topic{8.2\quad SRAM：写一个字节（$\sim$1--2 ns）}

场景：CPU 执行 \code{MOV [RSP+8], AL}，地址命中 L1D cache。

\begin{longtable}{L{0.06\linewidth}L{0.30\linewidth}L{0.30\linewidth}L{0.24\linewidth}}
\toprule
步 & 硬件事件 & 信号路径 & 耗时 \\
\midrule
1 & CPU 发出 store 地址+数据 & 地址+8 bit 数据 $\to$ L1D & 0 ns \\\tablerowrule
2 & Tag 比较命中 & 同读路径 & $\sim$0.3 ns \\\tablerowrule
3 & 行地址译码 & WL[42] 拉高 & $\sim$0.2 ns \\\tablerowrule
4 & 写驱动器驱动位线 & 8 个写驱动器把 BL 驱动为 AL 的各位值，$\overline{BL}$ 驱动为反值 & $\sim$0.1 ns \\\tablerowrule
5 & 单元翻转 & BL 通过 M5 把 Q 拉过反相器翻转阈值；反馈锁存到新稳态 & $\sim$0.5 ns \\\tablerowrule
6 & WL 关闭 & WL[42] 拉低，M5/M6 截止，新数据被锁存保持 & $\sim$0.1 ns \\\tablerowrule
7 & 写确认 & 写完成信号返回 CPU store buffer & $\sim$0.1 ns \\
\midrule
\multicolumn{3}{L{0.66\linewidth}}{总延迟（L1D 命中写）} & $\sim$1.3 ns \\
\bottomrule
\end{longtable}

关键区别：写操作中，外部写驱动器的电流必须\hwkey{强过}单元内部 PMOS 上拉管的反馈电流，才能把存储节点拉过翻转阈值。这就是为什么写操作对晶体管尺寸比（CR/PR）敏感，也是写辅助电路存在的原因。

\topic{8.3\quad DRAM：读一个字节（$\sim$50--100 ns）}

场景：CPU 执行 load，L1/L2/LLC 全部 miss，地址映射到 DDR5 DIMM 的某 bank 某行某列。

\begin{longtable}{L{0.06\linewidth}L{0.32\linewidth}L{0.30\linewidth}L{0.22\linewidth}}
\toprule
步 & 硬件事件 & 信号路径 & 耗时 \\
\midrule
1 & LLC miss，请求进入 MC 队列 & LLC $\to$ 内存控制器请求队列 & $\sim$2 ns \\\tablerowrule
2 & 调度器选中该请求 & FR-FCFS：检查目标 bank 状态 & $\sim$1 ns \\\tablerowrule
3 & 发 ACTIVATE 命令 & MC $\to$ PHY $\to$ CA 总线 $\to$ DRAM：打开目标行 & 0 ns（命令发出） \\\tablerowrule
4 & 行激活（tRCD） & DRAM 内部：行译码 $\to$ WL 打开 $\to$ 整行 16K 个电容与 BL 电荷共享 $\to$ 感放锁存到 row buffer & \hwtiming{16.25 ns}（DDR5-4800, tRCD=39CK） \\\tablerowrule
5 & 发 READ 命令 & MC $\to$ PHY $\to$ CA 总线：列地址 + burst 命令 & 0 ns \\\tablerowrule
6 & 列选择（CL） & DRAM 内部：列 mux 从 row buffer 选出目标列 $\to$ 输出缓冲；实际返回整个 burst（BL16 $\times$ 接口宽度 $\to$ 64 B 缓存行） & \hwtiming{16.67 ns}（DDR5-4800, CL=40CK） \\\tablerowrule
7 & DQS+DQ 传输 & DRAM 驱动 DQS 和 DQ $\to$ PCB 走线 $\to$ 控制器 PHY & $\sim$2 ns（走线延迟） \\\tablerowrule
8 & PHY 采样 & DQS 边沿采样 DQ，FIFO 对齐 & $\sim$1 ns \\\tablerowrule
9 & 数据返回 CPU & MC $\to$ LLC 填充 $\to$ L2 $\to$ L1 $\to$ 寄存器 & $\sim$5 ns \\
\midrule
\multicolumn{3}{L{0.66\linewidth}}{总延迟（DRAM 读，row miss）} & $\sim$45--55 ns \\
\bottomrule
\end{longtable}

\begin{pdfdiagram}[x=1cm,y=1cm]
  \node[hw state, minimum width=1.2cm] (cpu) at (0,0) {CPU};
  \node[hw state, minimum width=1.2cm] (llc) at (1.8,0) {LLC};
  \node[hw control, minimum width=1.4cm] (mc) at (3.8,0) {MC};
  \node[hw compute, minimum width=1.2cm] (phy) at (5.8,0) {PHY};
  \node[hw memory, minimum width=1.6cm] (dram) at (8.2,0) {DRAM};
  \node[hw state, minimum width=1.4cm] (cell) at (10.6,0) {1T1C};
  \draw[hw bus] (cpu) -- node[above,font=\tiny]{miss} (llc);
  \draw[hw bus] (llc) -- (mc);
  \draw[hw bus] (mc) -- node[above,font=\tiny]{ACT} (phy);
  \draw[hw bus] (phy) -- (dram);
  \draw[hw bus] (dram) -- node[above,font=\tiny]{WL+SA} (cell);
  \draw[hw return] (cell.south) -- (10.6,-0.9) -- node[above,font=\tiny]{DQS+DQ} (5.8,-0.9) -- (phy.south);
  \draw[hw return] (5.8,-0.9) -- (3.8,-0.9) -- (mc.south);
  \draw[hw return] (3.8,-0.9) -- (1.8,-0.9) -- (llc.south);
  \draw[hw return] (1.8,-0.9) -- (0,-0.9) -- (cpu.south);
  \node[hw label, text width=10cm] at (5.5,-1.6) {DRAM 读 1 字节：LLC miss→MC 调度→ACTIVATE（tRCD=16.25ns）→READ（CL=16.67ns）→DQS/DQ 传输→PHY 采样→返回。全程 $\sim$50 ns。};
\end{pdfdiagram}
\diagramnote{DRAM 读路径。延迟的大头在 DRAM 内部：行激活（电荷共享+感放）和列选择。PHY 传输和 MC 调度只占一小部分。}

若目标行已在 row buffer 中（row hit），步 4 可以跳过，总延迟降到 $\sim$25--30 ns（只有 CL + 传输）。这就是 FR-FCFS 调度优先发 row-hit 请求的原因。

\topic{8.4\quad DRAM：写一个字节（$\sim$50--100 ns）}

场景：同一 CPU 执行 \code{MOV [RSP+8], AL}，LLC miss，store 请求最终到达 DRAM。写路径与 8.3 的读路径共享行激活和列选择两个阶段，差异在于数据方向和收尾时序。

\begin{longtable}{L{0.06\linewidth}L{0.32\linewidth}L{0.30\linewidth}L{0.22\linewidth}}
\toprule
步 & 硬件事件 & 信号路径 & 耗时 \\
\midrule
1 & LLC miss，store 进入 MC 写队列 & 同读路径 & $\sim$2 ns \\\tablerowrule
2 & 调度器选中（可能需要读写切换） & 若当前总线方向为读，需等 turnaround（2--4 CK） & $\sim$1--2 ns \\\tablerowrule
3 & 发 ACTIVATE（若 row miss） & 同读路径步 3--4 & 16.25 ns \\\tablerowrule
4 & 发 WRITE 命令 + 数据 & MC $\to$ PHY 驱动 DQ/DQS $\to$ DRAM & $\sim$2 ns \\\tablerowrule
5 & DRAM 内部写入 & 列 mux 把 burst 数据从 DQ 写入 row buffer 对应位置 & $\sim$1 ns \\\tablerowrule
6 & 写恢复（tWR） & 数据经感放写入单元电容后，必须等待写恢复完成才能预充 & tWR=30 ns（DDR5-4800, 72CK） \\\tablerowrule
7 & 发 PRECHARGE（若需要） & 关闭当前行，位线恢复到预充电平 & tRP=16.25 ns（39CK） \\\tablerowrule
8 & 写确认 & MC 通知 LLC 写完成 & $\sim$1 ns \\
\midrule
\multicolumn{3}{L{0.66\linewidth}}{总延迟（DRAM 写，row miss，含 tWR/tRP）} & $\sim$65--70 ns \\
\bottomrule
\end{longtable}

DRAM 写的关键路径与读相近——都必须经过行激活和列选择；区别在于数据方向（读时 DRAM 驱动 DQ，写时控制器驱动 DQ），且写之后多一步写恢复（tWR），数据真正写入单元电容并稳定后才能预充。

\topic{8.5\quad NAND Flash：读一个字节（$\sim$25--100 $\mu$s）}

场景：CPU 执行 read() 系统调用，文件系统请求 SSD 读取 LBA 0x1234 处的 4 KiB 数据中的第 100 字节。

\begin{longtable}{L{0.06\linewidth}L{0.32\linewidth}L{0.30\linewidth}L{0.22\linewidth}}
\toprule
步 & 硬件事件 & 信号路径 & 耗时 \\
\midrule
1 & NVMe 命令到达 SSD 控制器 & CPU 写 SQ doorbell $\to$ SSD 控制器取命令 & $\sim$1 $\mu$s \\\tablerowrule
2 & FTL 查 L2P 映射 & LBA $\to$ 物理地址（die/plane/block/page） & $\sim$0.5 $\mu$s \\\tablerowrule
3 & 发 NAND 读命令 & 控制器 $\to$ NAND die：READ(page address) & $\sim$0.1 $\mu$s \\\tablerowrule
4 & NAND 内部读 & WL 逐层施加 Vread/Vread\_sel $\to$ 串电流导通/截止 $\to$ 感放读整页到 page register & \hwtiming{60--90 $\mu$s}（TLC） \\\tablerowrule
5 & 列选择 + 数据传输 & 从 page register 按列读出 $\to$ DQ 总线 $\to$ 控制器 SRAM & $\sim$5 $\mu$s（4 KiB） \\\tablerowrule
6 & LDPC ECC 校验 & 控制器对 4 KiB 数据做 LDPC 解码，纠正 bit 错误 & $\sim$2--5 $\mu$s \\\tablerowrule
7 & 数据返回 CPU & 控制器 DMA 到主机内存 $\to$ 完成队列 $\to$ MSI-X 中断 & $\sim$2 $\mu$s \\
\midrule
\multicolumn{3}{L{0.66\linewidth}}{总延迟（NAND 读 4 KiB）} & $\sim$70--100 $\mu$s \\
\bottomrule
\end{longtable}

\begin{pdfdiagram}[x=1cm,y=1cm]
  \node[hw state, minimum width=1.2cm] (cpu) at (0,0) {CPU};
  \node[hw control, minimum width=1.4cm] (nvme) at (2.0,0) {NVMe};
  \node[hw control, minimum width=1.4cm] (ftl) at (4.0,0) {FTL};
  \node[hw memory, minimum width=1.4cm] (nand) at (6.2,0) {NAND};
  \node[hw state, minimum width=1.4cm] (cell) at (8.4,0) {浮栅};
  \node[hw compute, minimum width=1.4cm] (ecc) at (10.6,0) {LDPC};
  \draw[hw bus] (cpu) -- node[above,font=\tiny]{SQ} (nvme);
  \draw[hw bus] (nvme) -- (ftl);
  \draw[hw bus] (ftl) -- node[above,font=\tiny]{READ cmd} (nand);
  \draw[hw bus] (nand) -- node[above,font=\tiny]{Vread} (cell);
  \draw[hw return] (cell.south) -- (8.4,-0.9) -- node[above,font=\tiny]{page data} (6.2,-0.9) -- (nand.south);
  \draw[hw bus] (nand) -- (ecc);
  \draw[hw return] (ecc.south) -- (10.6,-1.25) -- node[above,font=\tiny]{DMA} (0,-1.25) -- (cpu.south);
  \node[hw label, text width=10cm] at (5.5,-1.8) {NAND 读 1 字节（实际读整页 4--16 KiB）：NVMe 命令→FTL 映射→NAND 感放（60--90 $\mu$s）→LDPC 校验→DMA 返回。延迟比 DRAM 大 1000 倍。};
\end{pdfdiagram}
\diagramnote{NAND 读路径。延迟的绝对大头在 NAND 内部感放（步 4）：串电流需要时间稳定，感放需要时间分辨多电平。这就是为什么 SSD 随机读是 $\mu$s 级而不是 ns 级。}

注意：NAND 的最小读单位是 page（4--16 KiB），不能只读 1 字节。CPU 要的那 1 字节包含在整页数据中，由控制器从 page buffer 中按偏移取出。

步 4 为什么需要几十微秒？TLC 的 8 个电平要用 7 个读参考电压逐个区分，相邻电平间距仅 $\sim$0.8 V；每换一个参考电压，字线电压都要重新建立、串电流都要重新稳定，感放才能再比较一次。TLC 的页读时间（tR）因此达到 60--90 $\mu$s，是 SLC（只有 1 个读参考电压，约 25 $\mu$s）的数倍——这正是第五节单元类型表中“典型读延迟”一列的来源，也是本小节标题给出 25--100 $\mu$s 宽范围的原因：低端对应 SLC，高端对应 TLC/QLC。

\topic{8.6\quad NAND Flash：写一个字节（$\sim$200 $\mu$s--2 ms）}

场景：CPU 执行 write() 系统调用，4 KiB 数据经文件系统和 NVMe 队列到达 SSD。读一个字节时 NAND 只是多花些时间，写一个字节却会触发一连串结构性动作——根源是 NAND 不能原地改写。

\begin{longtable}{L{0.06\linewidth}L{0.32\linewidth}L{0.30\linewidth}L{0.22\linewidth}}
\toprule
步 & 硬件事件 & 信号路径 & 耗时 \\
\midrule
1 & NVMe 写命令到达 & 同读路径步 1--2 & $\sim$1.5 $\mu$s \\\tablerowrule
2 & FTL 分配空闲 page & 查空闲块表，选择目标 die/plane/block/page & $\sim$0.5 $\mu$s \\\tablerowrule
3 & 数据写入 page register & 控制器把 4--16 KiB 数据通过 DQ 写入 NAND page register & $\sim$5 $\mu$s \\\tablerowrule
4 & 发 PROGRAM 命令 & 控制器 $\to$ NAND：PROGRAM(page address) & $\sim$0.1 $\mu$s \\\tablerowrule
5 & NAND 内部编程（ISPP） & 逐次施加 Vpgm 脉冲（每次 +0.3V），每次后 verify，直到所有单元达到目标 Vth & \hwtiming{200 $\mu$s--2 ms}（TLC 典型 8--16 次脉冲） \\\tablerowrule
6 & 更新 L2P 映射 & 新物理地址写入映射表，旧 page 标记为无效 & $\sim$1 $\mu$s \\\tablerowrule
7 & 写确认 & 完成队列 + MSI-X & $\sim$1 $\mu$s \\
\midrule
\multicolumn{3}{L{0.66\linewidth}}{总延迟（NAND 写 4 KiB，TLC）} & $\sim$0.3--2 ms \\
\bottomrule
\end{longtable}

NAND 不能原地改写（in-place update）：浮栅中的电子只能通过擦除（整块 4--16 MiB 一起）才能移除。所以“写一个字节”实际上是：写到新的空闲 page $\to$ 更新映射 $\to$ 旧 page 等 GC 回收。这就是 FTL 和写放大存在的根本原因。

为什么 TLC 需要这么多次 ISPP 脉冲？8 个电平挤在几伏的阈值窗口内，相邻目标分布的间距只有 $\sim$0.8 V，分布稍宽就会跨进相邻电平。要把每个分布收紧，步进 $\Delta V_{ISPP}$ 必须取得小，并且每一步之后都要 verify，脉冲次数随之上升；SLC 只有 2 个电平、间距约 5 V，少数几个大步进脉冲即可完成。脉冲次数从 SLC 到 TLC 的倍增，正是编程时间从 200--300 $\mu$s 涨到 0.5--2 ms 的直接原因。

\topic{8.7\quad HDD：读一个字节（$\sim$4--12 ms）}

场景：CPU 执行 read() 系统调用，文件系统请求 HDD 读取 LBA 0x5678 处的 4 KiB 扇区中的第 200 字节。

\begin{longtable}{L{0.06\linewidth}L{0.32\linewidth}L{0.30\linewidth}L{0.22\linewidth}}
\toprule
步 & 硬件事件 & 信号路径 & 耗时 \\
\midrule
1 & SATA/SAS 命令到达 & 主机 $\to$ 磁盘控制器：READ DMA(LBA, count) & $\sim$1 $\mu$s \\\tablerowrule
2 & 地址转换 & LBA $\to$ 柱面/磁头/扇区（PBA） & $\sim$1 $\mu$s \\\tablerowrule
3 & 寻道（若磁头不在目标柱面） & VCM 加速 $\to$ 滑行 $\to$ 减速 $\to$ 伺服锁定目标磁道 & \hwtiming{4--9 ms}（平均） \\\tablerowrule
4 & 旋转等待（若目标扇区未转到磁头下方） & 盘片继续旋转，等待目标扇区到达 & \hwtiming{0--8.33 ms}（7200 rpm 平均 4.17 ms） \\\tablerowrule
5 & 磁头读取 & TMR 磁头飞越扇区 $\to$ 感阻变化 $\to$ 前置放大器 & $\sim$0.02 ms（4 KiB @ 200 MB/s） \\\tablerowrule
6 & 读通道处理 & AGC $\to$ 均衡 $\to$ Viterbi 检测 $\to$ LDPC 解码 & $\sim$0.05 ms \\\tablerowrule
7 & 数据返回 CPU & 扇区缓冲 $\to$ DMA $\to$ 主机内存 $\to$ 中断 & $\sim$0.01 ms \\
\midrule
\multicolumn{3}{L{0.66\linewidth}}{总延迟（HDD 随机读 4 KiB）} & \hwtiming{4--12 ms} \\
\bottomrule
\end{longtable}

\begin{pdfdiagram}[x=1cm,y=1cm]
  \node[hw state, minimum width=1.2cm] (cpu) at (0,0) {CPU};
  \node[hw control, minimum width=1.2cm] (sata) at (1.8,0) {SATA};
  \node[hw control, minimum width=1.4cm] (seek) at (3.8,0) {寻道};
  \node[hw state, minimum width=1.4cm] (rot) at (5.8,0) {旋转等待};
  \node[hw compute, minimum width=1.2cm] (head) at (7.8,0) {磁头};
  \node[hw compute, minimum width=1.4cm] (chan) at (9.8,0) {读通道};
  \node[hw state, minimum width=1.2cm] (buf) at (11.6,0) {缓冲};
  \draw[hw bus] (cpu) -- (sata);
  \draw[hw bus] (sata) -- (seek);
  \draw[hw bus] (seek) -- node[above,font=\tiny,yshift=3pt]{4--9ms} (rot);
  \draw[hw bus] (rot) -- node[above,font=\tiny,yshift=3pt]{0--8.33ms} (head);
  \draw[hw bus] (head) -- (chan);
  \draw[hw bus] (chan) -- (buf);
  \draw[hw return] (buf.south) -- (11.6,-0.9) -- node[above,font=\tiny]{DMA} (0,-0.9) -- (cpu.south);
  \node[hw label, text width=10cm] at (6.0,-1.4) {HDD 读 1 字节（实际读整扇区 4 KiB）：命令→寻道（4--9ms）→旋转等待（0--8.33ms，7200 RPM 平均 4.17ms）→磁头感读→读通道→DMA。机械运动占 99\% 延迟。};
\end{pdfdiagram}
\diagramnote{HDD 读路径。与 SRAM/DRAM/NAND 完全不同：延迟的 99\% 花在机械运动（寻道+旋转）上，电子信号处理只占不到 1\%。这就是为什么 HDD 随机 IOPS 只有 $\sim$100。}

若目标扇区恰好在磁头下方（顺序读），步 3 和步 4 都接近 0，总延迟降到 $\sim$0.02 ms（传输时间）。这就是顺序读 200 MB/s 而随机读只有约 0.5 MB/s 的原因。

两个口径问题需要说明。其一，步 2 沿用了经典的 CHS（cylinder-head-sector）寻址术语：\hwkey{柱面}（cylinder）指所有盘面上同一半径处的磁道集合——由于所有磁头固定在同一个 HSA 上同步径向移动，任一时刻可访问的磁道恰好构成一个圆柱面。其二，本表总延迟 4--12 ms 取的是“平均寻道 + 平均旋转延迟”的典型工作点（约 4.5 ms + 4.17 ms $\approx$ 8.7 ms，12 ms 对应较长的寻道）；6.19 的延迟分解给出的是完整包络 4--17 ms（最坏情况为全程寻道加一整转等待）。两种口径描述同一过程，前者代表典型随机读，后者代表取值范围。

\topic{8.8\quad HDD：写一个字节（$\sim$4--12 ms）}

场景：CPU 执行 write() 系统调用，4 KiB 数据到达 HDD 控制器。写路径的机械动作与 8.7 的读路径完全相同，差别只在最后一步的物理过程和掉电语义。

\begin{longtable}{L{0.06\linewidth}L{0.32\linewidth}L{0.30\linewidth}L{0.22\linewidth}}
\toprule
步 & 硬件事件 & 信号路径 & 耗时 \\
\midrule
1 & SATA 写命令 + 数据到达 & 主机 $\to$ 磁盘控制器 + 扇区缓冲 & $\sim$0.02 ms（4 KiB） \\\tablerowrule
2 & 地址转换 & 同读路径 & $\sim$1 $\mu$s \\\tablerowrule
3 & 寻道 & 同读路径 & 4--9 ms \\\tablerowrule
4 & 旋转等待 & 同读路径 & 0--8.33 ms \\\tablerowrule
5 & 写通道驱动磁头 & 写电流翻转磁头磁场 $\to$ 介质磁畴翻转 & $\sim$0.01 ms \\\tablerowrule
6 & 写确认 & 数据写入介质后（或写入缓存后）发完成 & $\sim$0.01 ms \\
\midrule
\multicolumn{3}{L{0.66\linewidth}}{总延迟（HDD 随机写 4 KiB）} & \hwtiming{4--12 ms} \\
\bottomrule
\end{longtable}

HDD 写和读的延迟几乎相同——瓶颈都是机械运动。区别在于：写时磁头产生磁场翻转介质磁畴，读时磁头检测磁畴产生的磁场。若开启写缓存（write cache），步 5--6 可以先写入 DRAM 缓存就返回确认，实际写入介质延后——但 HDD 没有 SSD 那样的掉电保护电容，突然掉电时缓存中尚未写入介质的数据即丢失（掉电瞬间驱动器只能靠主轴反电动势把磁头 ramp unload 到安全位置）。这正是主机提供 FUA（Force Unit Access）和 flush 命令强制数据写入介质的原因。

\begin{warningbox}
“写一个字节已返回成功”不等于“数据已在盘片上”。开启写缓存的 HDD 在数据进入片上 DRAM 缓存时就会报告完成，此后到真正写入介质之间存在毫秒级的窗口，掉电即丢失。对可靠性敏感的写入（日志、元数据）必须显式使用 FUA 位或 flush 命令——操作系统和数据库发出的 fsync，最终正是翻译成这些命令。
\end{warningbox}

\topic{8.9\quad 四条 trace 的对比}

把四条 trace 放在一起，延迟的量级差异一目了然。下表各行的延迟取典型工作点的代表值（DRAM 取 row miss，NAND 取 TLC 页读加控制器处理，HDD 取平均寻道加平均旋转延迟），用于跨器件的量级对比；完整的延迟分解和取值范围见 8.1--8.8 各小节（例如 NAND TLC 页读的典型 tR 为 60--90 $\mu$s，含控制器处理的总延迟约 70--100 $\mu$s，见 8.5）。

\begin{longtable}{L{0.12\linewidth}L{0.135\linewidth}L{0.135\linewidth}L{0.135\linewidth}L{0.135\linewidth}L{0.22\linewidth}}
\toprule
& SRAM & DRAM & NAND & HDD & 延迟瓶颈 \\
\midrule
读 1 字节 & $\sim$1.5 ns & $\sim$50 ns & $\sim$50 $\mu$s & $\sim$8 ms & SRAM: 门延迟；DRAM: 感放；NAND: 串电流；HDD: 机械 \\\tablerowrule
写 1 字节 & $\sim$1.3 ns & $\sim$65 ns & $\sim$1 ms & $\sim$8 ms & SRAM: 翻转阈值；DRAM: 读路径 + tWR；NAND: ISPP 多脉冲；HDD: 机械 \\\tablerowrule
最小访问粒度 & 1 字节 & 64 B（BL16 burst $\times$ 接口宽度 $\to$ 一条缓存行） & 4--16 KiB（page） & 4 KiB（扇区） & 粒度越大，随机访问浪费越多 \\\tablerowrule
写是否原地 & 是 & 是 & 否（需新 page） & 是（但 SMR 否） & NAND 的 FTL 由此而来 \\
\bottomrule
\end{longtable}

从 SRAM 到 HDD，延迟跨越了 \hwkey{7 个数量级}（1 ns $\to$ 10 ms）。每一级的延迟瓶颈完全不同：SRAM 是门电路传播延迟，DRAM 是电荷共享和感放时间，NAND 是隧穿编程的多脉冲验证，HDD 是机械臂的加速-滑行-减速。没有任何软件优化能消除这些物理极限——只能通过层次结构（cache、预取、队列）把慢设备的访问藏在快设备后面。

%======================================================================
\section{二、综合案例：一次 read() 从指令到盘片的完整旅程}
%======================================================================

前八节把每种器件的内部结构和各自的读写 trace 分别讲完了。本节用一个完整的案例把它们串起来：一个程序第一次读取某个文件的第一个字节。这一个动作会依次经过 CPU、cache 层次、内存控制器、DRAM、操作系统页缓存、NVMe 队列、SSD 控制器、FTL、NAND 阵列，再回到 CPU 寄存器——本章出现过的几乎所有概念都会在这条路径上用到一次。

\begin{keyidea}
这次旅程中 CPU 亲自参与的只有开头（发出系统调用）和结尾（把字节 load 进寄存器）各约 100 ns；中间约 100 $\mu$s 数据在队列、DMA 和 NAND 感放中流动，CPU 在运行别的进程。存储层次的本质在这一案例中体现得最完整：用每一级的小容量换时间，用异步和队列把等待藏起来。
\end{keyidea}

\topic{9.1\quad 案例设定：一个第一次读文件的程序}

\begin{lstlisting}[language=C]
int fd = open("data.bin", O_RDONLY);
char c;
read(fd, &c, 1);      // data.bin 第一次被读取,页缓存中没有它的数据
printf("%d\n", c);
\end{lstlisting}

假设这台机器配 DDR4-3200 内存和 TLC NVMe SSD，文件 \code{data.bin} 的第一个 4 KiB 页自开机以来从未被读过——这是这条路径能达到的“最冷”起点（更冷的情况只剩 HDD，见 9.7 的对比）。\code{read} 的第二个参数 \code{\&c} 是用户栈上的一个虚拟地址。

整个旅程分两大段：第一段由操作系统主导，把文件页从 SSD 搬进 DRAM 的页缓存（page cache，内核在 DRAM 中为文件数据维护的缓存，此后任何进程读同一页都不必再访问 SSD）；第二段由 CPU 硬件主导，把目标字节从 DRAM 经 cache 层次搬进寄存器。两段之间的交接点是 DMA 完成中断。

\begin{pdfdiagram}[x=0.96cm,y=1cm]
  \node[hw label, font=\small\bfseries, anchor=west] at (-0.4,5.05) {CPU / cache 通路（纳秒级）};
  \node[hw state, minimum width=1.7cm] (cpu) at (0,4.2) {CPU 核心};
  \node[hw state, minimum width=1.1cm] (l1) at (2.0,4.2) {L1};
  \node[hw state, minimum width=1.1cm] (l2) at (3.5,4.2) {L2};
  \node[hw state, minimum width=1.1cm] (l3) at (5.0,4.2) {L3};
  \node[hw control, minimum width=2.0cm] (mc) at (7.2,4.2) {内存控制器};
  \node[hw memory, minimum width=2.2cm] (dram) at (10.2,4.2) {DRAM 页缓存};
  \draw[hw bus] (cpu) -- (l1);
  \draw[hw bus] (l1) -- (l2);
  \draw[hw bus] (l2) -- (l3);
  \draw[hw bus] (l3) -- (mc);
  \draw[hw bus] (mc) -- (dram);
  \node[hw label] at (4.8,4.72) {① load 逐级未命中};
  \draw[hw return] (dram.south) -- (10.2,3.05) -- (0.35,3.05) -- (0.35,3.8);
  \node[hw label, fill=white, inner sep=2pt] at (5.3,3.05) {⑨ 64 B cache line：DRAM → L3 → L2 → L1 → 寄存器};

  % 分层标题位于两条左侧竖向通道和 DMA 回填通道之间。
  \node[hw label, font=\small\bfseries, anchor=west] at (1.25,1.62)
    {OS / NVMe / SSD 通路（微秒级）};
  \node[hw control, minimum width=2.1cm] (drv) at (0,0.65) {② NVMe 驱动};
  \node[hw control, minimum width=2.2cm] (sq) at (2.8,0.65) {③ SQ + doorbell};
  \node[hw compute, minimum width=2.1cm] (ssd) at (5.6,0.65) {④ SSD 控制器};
  \node[hw compute, minimum width=1.9cm] (ftl) at (8.2,0.65) {⑤ FTL（L2P）};
  \node[hw memory, minimum width=2.0cm] (nand) at (10.8,0.65) {⑥ NAND 页读};
  \draw[hw bus] (drv) -- (sq);
  \draw[hw bus] (sq) -- (ssd);
  \draw[hw bus] (ssd) -- (ftl);
  \draw[hw bus] (ftl) -- (nand);

  % 系统调用与返回分别占据左侧两条竖向通道。
  \draw[hw bus] ($(cpu.south)+(-0.35,0)$) -- (-0.35,1.1);
  \node[hw label, rotate=90] at (-0.72,2.42) {read() 系统调用};
  \draw[hw return] ($(drv.north)+(0.35,0)$) -- (0.35,3.8);
  \node[hw label, rotate=90] at (0.75,2.30) {⑧ 中断后返回};

  % DMA 回填在两层中间，完成通知放到底部，互不交叉。
  \draw[hw return] (ssd.north) -- (5.6,1.18) -- (6.95,1.18)
    -- (6.95,2.15) -- (11.75,2.15) -- (11.75,4.2) -- (dram.east);
  \node[hw label, fill=white, inner sep=2pt] at (8.55,2.15) {⑦ DMA 写 DRAM（16 KiB）};
  \draw[hw return] (ssd.south) -- (5.6,-0.55) -- (0,-0.55) -- (drv.south);
  \node[hw label, fill=white, inner sep=2pt] at (2.8,-0.55) {CQE + MSI-X：唤醒驱动};

  \node[hw label, text width=12cm] at (5.4,-1.65) {完整旅程：read() → 页缓存未命中 → NVMe 队列 → FTL → NAND 页读（60--90 $\mu$s）→ DMA 回填 DRAM → 中断 → copy\_to\_user → 返回用户态 → load 经 cache 层次把字节搬进寄存器。};
\end{pdfdiagram}
\diagramnote{从 read() 到寄存器的完整路径。上排是 CPU 硬件主导的 cache 层次(纳秒级),下排是操作系统与 SSD 主导的设备路径(微秒级),两段之间由 DMA 和中断衔接。两段的延迟相差约 1000 倍。}

\topic{9.2\quad 第一段:从 syscall 指令到 NVMe 命令}

\code{read} 库函数把参数按系统调用约定装入寄存器后，用一条指令陷入内核（x86-64, Intel 语法）：

\begin{lstlisting}
mov   eax, 0          ; SYS_read;rdi = fd,rsi = buf,rdx = 1 已先行设置
syscall               ; 陷入内核:RIP 跳到系统调用入口,切换到内核栈
\end{lstlisting}

进入内核后的前几个动作都在纳秒到微秒量级完成：系统调用入口保存用户态现场、查系统调用表转到 \code{sys\_read}（约 0.1 $\mu$s）；VFS 层按 \code{fd} 找到文件的 inode，在页缓存的基数树中查找该页——未命中（约 0.2 $\mu$s）；块层把请求交给 NVMe 驱动，驱动在 DRAM 的提交队列（SQ）中构造一个 64 字节的读命令项（含 PRP——指向目标物理页的物理地址列表），然后做一步关键动作：向 SSD 控制器的 doorbell 寄存器做一次 \term{MMIO 写}{memory-mapped I/O write}。doorbell 是 PCIe BAR 空间映射到 CPU 物理地址空间的设备寄存器，这次写不会进 cache，而是经内存控制器路由到 PCIe 根复合体，以 TLP 包的形式到达 SSD——这是 CPU 与 SSD 之间唯一的“通知”手段，总共约 0.5 $\mu$s。

此后 CPU 对这次读的所有权就结束了：发起 read 的进程被挂起（睡眠在 NVMe 完成等待上），调度器把 CPU 核心交给其他就绪进程。在 NAND 感放的几十微秒里，这个核心可以执行别的程序的几十万条指令。

\topic{9.3\quad 第一段后半:SSD 内部与 DMA 回填}

SSD 控制器看到 doorbell 后，从 DRAM 的 SQ 中 DMA 取走命令（它不通过 CPU 读内存，而是自己发起 PCIe 读事务），FTL 把命令中的逻辑页号查 L2P 映射表翻译成物理位置（die/plane/block/page），向目标 NAND die 发 READ 命令。接下来是全程最贵的一步：NAND 页读，TLC 典型 \hwtiming{60--90 $\mu$s}——8 个电平要用 7 个读参考电压逐个区分，每换一次参考电压字线都要重新建立（机制见 8.5）。整页（如 16 KiB）读出后经 LDPC 校验放入控制器缓冲，FTL 更新必要的元数据。

随后 SSD 按命令里的 PRP 列表发起 DMA 写事务，把 16 KiB 页数据经 PCIe 直接写入 DRAM 中页缓存的目标页帧（约 3 $\mu$s）——这一步不消耗任何 CPU 指令。写完，SSD 在 DRAM 的完成队列（CQ）中放入完成项，并触发 \term{MSI-X 中断}{message-signaled interrupt}：一个写到 CPU 本地 APIC 地址的 PCIe 消息，把对应核心从当前进程踢进内核的中断处理例程。ISR 标记请求完成、唤醒睡眠的进程（约 1.5 $\mu$s），内核把用户要的那 1 个字节从页缓存复制到 \code{\&c}（copy\_to\_user，约 0.1 $\mu$s），系统调用返回。

\topic{9.4\quad 第二段:地址拆解与命中判定——硬件怎么知道“不在 cache 里”}

回到用户态后，\code{printf} 要读变量 \code{c}，编译生成的 load 指令开始走硬件路径：

\begin{lstlisting}
mov   al, [rsi]       ; 读回刚到位的字节:虚拟地址 rsi
\end{lstlisting}

第一步是\term{地址翻译}{address translation}：虚拟地址送 TLB（地址翻译缓存），命中（约 1 个时钟周期）直接得到物理页帧号；若 TLB 未命中，MMU 硬件走 4 级页表（page walk），每级一次内存访问、共约 20--80 ns。本案例中用户栈的映射早已建立，TLB 命中是常态。

拿到物理地址后，cache 控制器做的第一件事是把地址\hwkey{拆成三段}。以一颗典型 x86-64 核的 32 KiB、8 路组相联 L1D 为例：

\begin{longtable}{L{0.16\linewidth}L{0.16\linewidth}L{0.56\linewidth}}
\toprule
字段 & 位段 & 作用 \\
\midrule
偏移 offset & PA[5:0] & 64 B 缓存行内的字节位置（$2^6=64$） \\\tablerowrule
索引 index & PA[11:6] & 选哪一个组（set）：$32\,\text{KiB} \div 8\,\text{way} \div 64\,\text{B} = 64$ 组，$2^6=64$ \\\tablerowrule
标记 tag & PA[47:12] & 与组内 8 个 way 存的 tag 逐一比较，判断“是不是我要的那一行” \\
\bottomrule
\end{longtable}

这个划分有一个精心设计的地方：index 和 offset 合起来恰好 12 bit，不超出 4 KiB 页内偏移——页内偏移在虚拟地址到物理地址的翻译中\emph{不变}，因此 L1 可以用虚拟地址的低位去索引组、同时让 TLB 翻译高位，两者\hwkey{并行}进行，谁也不等谁。这就是 \term{VIPT}{virtually indexed, physically tagged}（虚拟索引、物理标记），也是“L1 容量 $\div$ 相联度 $\le$ 页大小”这条经典约束的来源。

判命中是完全由组合逻辑和 SRAM 阵列配合完成的，不需要任何“查询指令”：

\begin{pdfdiagram}[x=1cm,y=1cm]
  % 地址段标签
  \node[hw label, font=\tiny] (idx) at (0.5,2.3) {PA[11:6] index};
  \node[hw label, font=\tiny] (tagin) at (3.9,2.3) {PA[47:12] tag};
  \node[hw label, font=\tiny] (offin) at (8.3,0.3) {PA[5:0] offset};
  % 上排:tag 路径
  \node[hw memory, minimum width=2.6cm] (tagarr) at (0.5,1.2) {tag+valid 阵列};
  \node[hw compute, minimum width=1.8cm] (cmp) at (3.9,1.2) {比较器 $\times$8};
  \node[hw control, minimum width=1.4cm] (hit) at (6.4,1.2) {hit 向量};
  \draw[hw bus] (tagarr) -- (cmp);
  \draw[hw bus] (cmp) -- (hit);
  \draw[hw bus] (tagin.south) -- (cmp.north);
  % index 通道:左侧竖直干线同时进 tag 阵列与数据阵列
  \draw[hw bus] (idx.south) -- (tagarr.north);
  \draw[hw bus] (idx.west) -- ++(-1.4,0) -- ++(0,-3.1) -- (-1.1,-0.8);
  \node[hw label, font=\tiny, rotate=90] at (-1.15,0.2) {选 set};
  % 下排:数据路径
  \node[hw memory, minimum width=2.6cm] (datarr) at (0.5,-0.8) {数据阵列(8 way)};
  \node[hw compute, minimum width=1.6cm] (waymux) at (3.9,-0.8) {way mux};
  \node[hw compute, minimum width=2.0cm] (bytemux) at (6.9,-0.8) {64:1 字节 mux};
  \node[hw state, minimum width=1.6cm] (reg) at (9.6,-0.8) {寄存器堆};
  \draw[hw bus] (datarr) -- (waymux);
  \draw[hw bus] (waymux) -- node[above,font=\tiny]{64 B 行} (bytemux);
  \draw[hw bus] (bytemux) -- node[above,font=\tiny]{1 B} (reg);
  % hit 向量从顶部通道进 waymux.north
  \draw[hw return] (hit.south) -- (6.4,0.1) -- (3.9,0.1) -- (waymux.north);
  \node[hw label, font=\tiny] at (5.15,0.28) {命中 way 选择};
  % offset 进 bytemux.north
  \draw[hw bus] (offin.south) -- (8.3,-0.2) -- (6.9,-0.2) -- (bytemux.north);
  \node[hw label, text width=11cm] at (4.5,-2.1) {命中判定电路。index 同时索引 tag 阵列与数据阵列：8 个 way 的 tag 与数据行\emph{同时}读出，8 个比较器并行比对 PA tag 与 valid 位，得到 hit 向量；数据行早已候在 way mux 输入端，hit 向量只需选出其中一路，再由 offset 经两级 mux(先 8:1 选 8 字节字、再 8:1 选字节)取出目标字节送寄存器堆写口。};
\end{pdfdiagram}
\diagramnote{L1 命中判定的数据通路。关键路径是“tag 比较 + way mux”，而不是数据读出——8 个 way 的 tag 行和数据行是被同一条 index 译码线同时驱动的，比较结果到达时数据已经等在 mux 输入端。这就是 8 路组相联仍能做到约 1 ns 命中的原因。}

逐拍看这条通路：index 经行译码器驱动 tag 阵列和数据阵列各一条字线（就是 1.8 节外围电路里的行译码），8 个 way 的“tag + valid + 数据行”同时读出；8 组异或-或非比较器把读出的 tag 与 PA tag 并行比对，再与各自的 valid 位相与，经或树得到\hwkey{命中/未命中}信号和命中 way 编号。若命中，way mux 选出对应 64 B 行，offset[5:3] 先在 8 个 8 字节字里选一个字，offset[2:0] 再在 8 个字节里选目标字节——两级 8:1 mux 比一级 64:1 mux 的扇入小、延迟低，选中的字节随下一个时钟上升沿写入寄存器堆的 \code{al} 写口。若 8 路全不命中（或命中行 valid=0），控制器拉高 miss 信号，转入 9.5 的流程。

\topic{9.5\quad miss 之后:MSHR、逐级请求与回填控制}

miss 由谁“控制”？答案是一组专用的硬件记账结构，而不是流水线停止等待。检测到 miss 后，L1 控制器把这次访问登记进一个 \term{MSHR}{miss status holding register}（未命中状态保持寄存器）：表项里记录请求地址、目标寄存器、访问类型和字节掩码。登记完流水线就解放了——后续指令若命中 cache 可以继续执行（\hwkey{hit under miss}），只有用到这个 miss 结果的指令才会在退休前等待。MSHR 的项数（典型 10--16 项）就是一颗核同时允许的 outstanding miss 数，也是内存级并行度的上限。

请求随后逐级下探，每一级重复同样的判定：L1 的 MSHR 向 L2 发请求，L2 查自己的 tag 阵列，命中则返回整个 64 B 行，未命中则登记自己的 MSHR 再向 L3；L3 未命中才把请求放上片内互连（ring/mesh）送到内存控制器。内存控制器侧的动作就是 8.3 的完整 trace：本例中目标页帧刚被 DMA 写过，按 row miss 处理——ACT（tRCD 13.75 ns）打开所在 bank 的行，整行经感放锁存进 row buffer，READ + CL（13.75 ns）选出目标列，BL8 突发经 DQS 把 64 B 送回 PHY（共约 50--100 ns）。

回填路径上还有三个控制细节。其一，\term{关键字优先}{critical word first}：DRAM 按突发返回 64 B，控制器把 CPU 等的那个 8 字节字排在突发最前面，到达后立即经旁路送给等待的 load，其余部分继续填充，不必等整行到齐。其二，回填写入 L1 数据阵列的同时更新 tag 阵列（写入 PA tag、置 valid）；若选中的牺牲行是脏的（dirty=1，之前被写过），它不能简单丢弃，要先复制到\term{写回缓冲}{writeback buffer}，稍后再写入 L2——回填与写回因此可以并行。其三，每一级的 MSHR 在数据返回后还要完成“对账”：把行装入数据阵列、按表项里的目标寄存器唤醒等待的 load、释放表项。整条 miss 路径没有一条软件指令参与，全部是各 cache 控制器里的状态机按固定协议接力完成。

\topic{9.6\quad 选中与修改那一个字节:字节 mux、字节使能与写路径}

读一个字节的最后一段电路在 9.4 已经看到：64 B 行经两级 mux 选出目标字节。值得强调的是“选中”的物理含义——offset 的 6 个 bit 就是两级 mux 的选择信号，每一级都是 1.8 节讲过的传输门 mux 树：offset[5:3] 三根线控制 8 选 1（8 字节字），offset[2:0] 三根线再控制 8 选 1（字节）。被选中的字节经驱动管送上内部总线，到达寄存器堆写口时伴随着写使能和目的寄存器号；\code{al} 是 \code{rax} 的低 8 位，写口只更新这 8 个触发器，其余 56 位的写使能保持关闭。

如果这一步不是读而是\hwkey{写}（\code{mov [rsi], al}），电路要回答三个问题。第一，写到哪里：写命中时（写回式 cache），控制器用 tag 比较定位命中行后，并不改写整行——\term{字节使能}{byte enable} 信号（64 根，每字节一根）只有 offset 选中的那根有效，写入电路（1.4 节的写驱动器）只把对应字节的 8 对位线差分强拉到新值，字线打开后访问管与上拉管的竞争把锁存器翻转成新数据；同行其余 63 字节的位线保持预充，数据不动。写命中同时把该行的 dirty 位置 1——这行从此与下级不一致，将来被替换时必须写回。第二，写未命中怎么办：x86 的 L1 采用 \term{write allocate}{写分配}——先按读 miss 把整行取进 cache，再按写命中处理；另一种策略 no-write-allocate 则绕过 cache 直接写下级，常用于写后不再读的数据。第三，顺序怎么保证：store 并不在发射时写 cache，而是先进 \term{store buffer}{存储缓冲}，等指令退休（成为全局可见顺序中的最旧未完结指令）后才把字节写入数据阵列——这是推测执行下“写不能被撤销”这条规则的硬件实现。

还有一个容易忽略的代价：带 SECDED ECC 的 cache（服务器常见）里，改一个字节后整行的校验位就失效了，控制器必须按新行内容重新计算 ECC——一次字节写实际上触发“读出行、修改字节、重算 ECC、写回行”的完整序列，这也是为什么 ECC 行的写开销略高于普通行。

注意这条路径上的粒度放大：程序要 1 个字节，CPU 与 DRAM 之间搬的是 64 B 缓存行，SSD 与 DRAM 之间搬的是 16 KiB 页，NAND 内部读的是一整页、未来 GC 搬的是整块。每一层都用自己的最小粒度为上一层服务——这就是 8.9 表中“最小访问粒度”一列在真实系统中的样子。

\topic{9.7\quad 时间总账与同一行代码的六种结局}

把两段加起来，第一次读这个文件的总耗时：

\tablechunkintro{1}{2}{1}{7}
\begin{longtable}{L{0.05\linewidth}L{0.15\linewidth}L{0.42\linewidth}L{0.13\linewidth}L{0.13\linewidth}}
\toprule
步 & 主导者 & 事件 & 耗时 & 累计 \\
\midrule
1 & CPU$\to$OS & \code{syscall} 陷入内核，查系统调用表 & $\sim$0.1 $\mu$s & 0.1 $\mu$s \\\tablerowrule
2 & OS & VFS $\to$ 页缓存查找，未命中 & $\sim$0.2 $\mu$s & 0.3 $\mu$s \\\tablerowrule
3 & OS$\to$SSD & 构造 SQ 项，MMIO 写 doorbell & $\sim$0.5 $\mu$s & 0.8 $\mu$s \\\tablerowrule
4 & SSD & 取命令，FTL 查 L2P 映射 & $\sim$1.5 $\mu$s & 2.3 $\mu$s \\\tablerowrule
5 & SSD & NAND 页读（TLC tR，串电流 + 多参考电压感放） & \hwtiming{60--90 $\mu$s} & $\sim$92 $\mu$s \\\tablerowrule
6 & SSD & LDPC 校验，数据入控制器缓冲 & $\sim$1 $\mu$s & $\sim$93 $\mu$s \\\tablerowrule
7 & SSD$\to$DRAM & DMA 把 16 KiB 页写入页缓存 & $\sim$3 $\mu$s & $\sim$96 $\mu$s \\
\bottomrule
\end{longtable}

\tablechunkintro{2}{2}{8}{汇总行}
\begin{longtable}{L{0.05\linewidth}L{0.15\linewidth}L{0.42\linewidth}L{0.13\linewidth}L{0.13\linewidth}}
\toprule
步 & 主导者 & 事件 & 耗时 & 累计 \\
\midrule
8 & SSD$\to$OS & CQE + MSI-X 中断，ISR 唤醒进程 & $\sim$1.5 $\mu$s & $\sim$98 $\mu$s \\\tablerowrule
9 & OS & copy\_to\_user：1 字节复制到用户缓冲 & $\sim$0.1 $\mu$s & $\sim$98 $\mu$s \\\tablerowrule
10 & CPU & 返回用户态；load：L1/L2/L3 未命中 $\to$ DRAM $\to$ 寄存器 & $\sim$0.1 $\mu$s & $\sim$98 $\mu$s \\\tablerowrule
\multicolumn{4}{L{0.75\linewidth}}{总延迟（第一次读文件的一个字节，TLC SSD）} & \hwtiming{$\sim$70--100 $\mu$s} \\
\bottomrule
\end{longtable}

总账的结构值得看一眼：步 5 一步占了约 90\%，OS 与 PCIe 上的所有动作（步 1--4、6--10）合计不到 10 $\mu$s。如果同一进程紧接着读这个文件的第二个字节，路径完全不同——页缓存已命中，只剩步 1、2、9、10，共约 1--2 $\mu$s；如果该缓存行还在 L1 里，只剩步 10 的 L1 命中，约 1 ns。同一个 \code{read}、同一条 load，结局横跨七个数量级：

\begin{longtable}{L{0.24\linewidth}L{0.44\linewidth}L{0.20\linewidth}}
\toprule
数据在哪里 & 实际路径 & 延迟 \\
\midrule
L1 & load 直接命中 & $\sim$1 ns \\\tablerowrule
L2 / L3 & L1 未命中，下级命中 & $\sim$5--40 ns \\\tablerowrule
DRAM（row miss） & 内存控制器 $\to$ tRCD + CL $\to$ burst & $\sim$50--100 ns \\\tablerowrule
页缓存（文件页已在内存） & syscall + VFS 命中 + copy & $\sim$1--2 $\mu$s \\\tablerowrule
SSD（页缓存未命中） & 本案例全程 & $\sim$70--100 $\mu$s \\\tablerowrule
HDD（页缓存未命中） & 同上，但 NAND 换成寻道 + 旋转 & $\sim$5--13 ms \\
\bottomrule
\end{longtable}

这张表就是全章的收尾。从 6T 锁存器到浮栅到磁畴，每种器件的物理结构决定了它在阶梯上的位置；cache、页缓存、队列、DMA 这些机制不改变任何一级的物理延迟，只决定一次访问实际停在哪一级。程序员的 \code{read()} 只有一行，硬件和操作系统在背后完成的，是本章从一颗电容讲到一整张表的全部内容。
\section{三、块设备与文件系统的硬件接口}

本节回答一个朴素的问题：一次写文件，到了设备那端究竟发生
了什么。答案分两半。一半是编址：块设备从不按字节寻址，
主机看到的是一组从 0 起编的扇区，这套编址从 CHS 几何参
数一路演化到 LBA；另一半是缓存：从 CPU 到介质隔着页缓
存、设备缓存两层 RAM，“写完了”在哪一层被宣布，直接决定
断电后数据还在不在。

\topic{块设备与字节设备的分界}
操作系统把外设粗分成两类。字节设备（character device，
通行译名“字符设备”）
按流读写：串口每来一个中断交出一个字节，键盘把按键编码
逐个上交，终端把输出字节按顺序显示。流没有“位置”可言：
读过的字节不会再来，写出去的字节收不回，“回到第 37 个
字节”这种请求在流上不成立。块设备（block device）按固
定大小的整块读写：磁盘、SSD、eMMC 都以块为最小传输单位，
每块有编号，可按号随机读、随机写、反复改写。前者像自来
水管，后者像一墙编了号的储物柜；\code{/dev} 里设备文件
类型位的 \code{b} 与 \code{c}，标的就是这个分界。

分界不是软件随意划的，它几乎由介质单方面决定。存储介质
普遍成块，因为定位与寻址的代价按次数收、不按字节收：机
械盘的磁头移一次、闪存的页选通一次、磁带走带定位一次，
都是固定开销，摊到越多字节上越划算；介质于是把字节捆成
块、给块编号，让一次定位服务一整块。终端和串口相反：它
们的“介质”是线路另一端的人或进程，数据以事件节奏零散到
达，没有可寻址的静止形态，天然成流。设备分类表因此大体
是一张介质决定论的清单：

\begin{longtable}{L{0.13\linewidth}L{0.25\linewidth}L{0.23\linewidth}L{0.27\linewidth}}
\toprule
类别 & 代表设备 & 读写单位 & 为什么是这样 \\
\midrule
块设备 & 机械盘、SSD、eMMC、U 盘 & 整块，按号寻址，可改写 & 介质按块组织，定位代价按次收 \\\tablerowrule
字节设备 & 串口、键盘、终端、打印机 & 字节流，无地址，不可回看 & 数据以事件节奏到达，无静止形态 \\\tablerowrule
网络设备 & 网卡 & 帧（分组），突发到达 & 既非流也非块，第四节专门讲 \\\tablerowrule
特例 & 内存条 & 按地址随机访问 & 它就是“地址”本身，不需块这层抽象 \\
\bottomrule
\end{longtable}

两处边界值得留神。其一，两种接口可以互相伪装，只是代价
不对等：块设备装成流很容易——读方向是按号逐块读出、
把读出的字节串成流，写方向是把流攒成块再逐块下发，
\code{dd} 按字节计数拷盘靠的就是这层软件转换；字节设
备装成块则要在流
上模拟随机访问，理论上可行，工程上没人做。磁带是骑在边
界上的活例子：它按记录块寻址、算块设备，却只支持顺序定
位，跳读远处的块以十秒计，于是当不了文件系统的载体，只
配做归档——“可随机定位”才是块设备配享文件系统的入场券。
其二，“可改写”对块设备是逻辑承诺、不是物理事实：主机以
为自己在原地改写了 37 号块，SSD 的 FTL 完全可能把新数
据写到别处、只改一张映射表。这是“SSD 与闪存的硬件特性”一节的主题，这里先记
住结论：块设备承诺的是“按号寻址的整块读写”，不承诺任何
物理位置。

还要防止一个误读：流设备也有缓冲，缓冲不改变流的本质。
串口的 UART 自带 16 字节 FIFO，终端有行缓冲，管道在内核
里有 64 KiB 环形缓冲——它们吸收的是速率抖动，不是引入
地址：缓冲里的字节依然先到先走、读完即弃，没有任何“回
到第 37 字节”的办法。判断一个设备是流还是块，不问它有
没有缓冲，只问它的数据有没有编号、能不能按编号回头取。

\topic{扇区、块与 LBA}
块设备的最小寻址单位叫扇区（sector），称呼来自机械盘：
盘片上每个磁道被切成若干弧段，一段就是一个扇区。早期接
口把物理几何直接暴露给软件——经典芯片与平台案例 int 13h 的调用约定
就是标本：柱面号、磁头号、扇区号分装三个寄存器，柱面
10 位、磁头 8 位、扇区 6 位且从 1 起编。这套 CHS
（cylinder-head-sector）编址把容量上限锁死在寄存器位宽
里：

\begin{lstlisting}
CHS 上限 = 1024 柱面 x 256 磁头 x 63 扇区 x 512 字节
         = 8,455,716,864 字节，约 8.4 GB
DOS 兼容约定把磁头数限到 255，上限收缩到约 7.8 GiB
\end{lstlisting}

更早还有一道更矮的墙：BIOS 与 IDE 寄存器各自的几何位宽
取交集，最窄处拼出的上限更小：

\begin{longtable}{L{0.26\linewidth}L{0.16\linewidth}L{0.16\linewidth}L{0.2\linewidth}}
\toprule
接口 & 柱面位数 & 磁头位数 & 扇区位数 \\
\midrule
BIOS int 13h & 10 & 8 & 6（从 1 起编） \\\tablerowrule
IDE/ATA & 16 & 4 & 8 \\\tablerowrule
两者交集（实际可用） & 10 & 4 & 6 \\
\bottomrule
\end{longtable}

交集那行就是 $1024\times16\times63\times512$ 字节
$\approx504$ MiB 这道墙的出处，1990 年代中期的盘纷纷撞
线。解法不是继续加位宽，而是换掉抽象本身：LBA（logical
block addressing）把盘看成从 0 起编的一维扇区数组，几何
换算藏进盘内控制器。int 13h 扩展（\code{AH=0x42}）改用
磁盘地址包传 64 位 LBA；ATA 协议的 LBA28 给 28 位地址，
上限 $2^{28}\times512$ 字节 $=128$ GiB；LBA48 给 48 位，
上限 $2^{57}$ 字节 $=128$ PiB，容量焦虑就此退场。CHS 从
此退化为兼容字段：现代盘报出的柱面磁头数是编出来哄老软
件的，与盘片真实布局毫无关系。经典芯片与平台案例讲固件把硬件细节封
装在寄存器约定后面，LBA 是同一方法用在编址上——把“盘面
几何”换成“线性数组”，接口从此与介质解耦。

扇区本身多大，三十年里的答案是 512 字节；2010 年前后
Advanced Format 把物理扇区扩到 4 KiB：同样的盘片面积能
放下更多数据（每扇区的同步头、间隙与纠错码占比随扇区变
大而变小），纠错能力也更强。麻烦在操作系统生态仍按 512
字节发命令，过渡期于是生出两种盘：

\begin{longtable}{L{0.1\linewidth}L{0.24\linewidth}L{0.52\linewidth}}
\toprule
盘类型 & 逻辑/物理扇区 & 部分写的代价 \\
\midrule
512n & 512 B / 512 B & 无翻译，老盘直接写 \\\tablerowrule
512e & 512 B / 4 KiB & 固件翻译；写不满一个物理扇区时盘内读-改-写 \\\tablerowrule
4Kn & 4 KiB / 4 KiB & 无翻译；只认 512 B 的老软件直接出错 \\
\bottomrule
\end{longtable}

512e 引出了“原子单位”这个关键概念。块设备的读写以逻辑
扇区为原子单位：一次写要么整个扇区成功、要么整个扇区失
败，不存在“半个新扇区加半个旧扇区”的可见中间态（断电恰
好发生在两个扇区之间是另一回事，由文件系统日志保证一致性）。
主机写一个 512 字节逻辑扇区、物理扇区却是 4 KiB 时，那
4 KiB 里有八分之七是必须保留的老数据，盘内只能读-改-
写：先把物理扇区整页读进设备缓存，替换其中 512 字节，
再整页写回。一次逻辑写放大成“一次读加一次写”；在
7200 rpm 的机械盘上若错过当圈的写窗口，还要多等一整圈
$60/7200\approx8.3$ ms。读-改-写（read-modify-write）
这个词后面还会反复出现：文件系统未对齐写、RAID 的部分
条带写、闪存的页内改写，全是同一机制在不同层级的化身。

还要记住 LBA 的编号单位是逻辑扇区、不是字节：同一个
LBA 2048，在 512e 盘上是字节偏移 1 MiB，在 4Kn 盘上是
8 MiB。分区表、引导代码、\code{dd} 的 \code{skip} 计数，
全都按逻辑扇区计——换盘或换镜像格式时，这是最容易算错
的一位。

\topic{文件系统块与设备扇区}
文件系统在块设备之上又定了一级自己的单位：文件系统块，
Linux 桌面与服务器最常见的取值是 4 KiB。文件的分配、空
闲位图、日志全按这个粒度登记，一次写文件最终化成若干个
4 KiB 块下发到块层。两层单位同名不同级：设备扇区是硬件
的原子单位，文件系统块是软件的分配单位。两者对齐则相安
无事，错开则每次写都在交学费。

对齐是道算术题。设物理扇区 4 KiB，分区从 63 号 512 字
节扇区开始——MBR 时代按“柱面对齐”留下的经典起点：
$63\times512=32\,256$ 字节，除以 4096 得 7.875，不是整
数。于是分区里每个 4 KiB 文件系统块都横跨两个物理扇区：
写这样一个块要动两个物理扇区，在 512e 盘上每个被部分覆
盖的物理扇区都要读-改-写；一次文件系统写拆成两笔、每
笔先读后写，随机写吞吐近乎腰斩，顺序写也因多出的盘内操
作而打折。闪存上还有第二重错位：擦除块以几百 KiB 计
（“SSD 与闪存的硬件特性”一节细讲），分区起点与擦除块边界错开，垃圾回收每轮
都要多搬一批页。

解决办法笨拙而有效：对齐到一个足够大的公倍数。1 MiB
$=2^{20}$ 字节，是 512 字节扇区的 2048 倍、4 KiB 扇区的
256 倍，也能整除常见 RAID 条带（64/128/256 KiB）和多数
闪存擦除块（256 KiB-1 MiB）。2000 年代末起，主流分区工
具把默认起点从 63 号扇区改到 2048 号扇区（整 1 MiB），
“分区对齐 1 MiB”从此成为惯例。三种起点的对照一目了然：

\begin{longtable}{L{0.24\linewidth}L{0.16\linewidth}L{0.2\linewidth}L{0.26\linewidth}}
\toprule
分区起点 & 起点字节 & 整除 4 KiB？ & 后果 \\
\midrule
63 号 512 B 扇区 & 32 256 B & 否（7.875） & 每个文件系统块横跨两个物理扇区 \\\tablerowrule
64 号 512 B 扇区 & 32 768 B & 能（8） & 对齐扇区，但不对齐 1 MiB \\\tablerowrule
2048 号 512 B 扇区 & 1 MiB & 能（256） & 扇区、条带、擦除块全部对齐 \\
\bottomrule
\end{longtable}

检查办法只剩一句算术：分区起始 LBA 乘扇区大小，能同时被
4096、条带大小与擦除块大小整除，才算对齐；\code{fdisk -l}
与 \code{lsblk} 都能直接读出起点，装完系统第一件事就该
验这道题。

\topic{读写缓存的层级}
从 CPU 到介质，数据要经过两层 RAM，每一层都可能“代收”
一次写：

\begin{longtable}{L{0.14\linewidth}L{0.26\linewidth}L{0.2\linewidth}L{0.24\linewidth}}
\toprule
层级 & 位置与容量级 & 读命中时延量级 & 断电后 \\
\midrule
页缓存 & 主存，GiB 级 & 约 100 ns & 丢失 \\\tablerowrule
设备缓存 & 盘内 DRAM，MiB-GiB 级 & 数十 $\mu$s & 无掉电保护则丢失 \\\tablerowrule
介质 & 闪存或盘片，百 GiB-TiB 级 & $10^{2}$ $\mu$s-$10$ ms & 保留 \\
\bottomrule
\end{longtable}

页缓存（page cache）是操作系统拿空闲主存开的代收点：写
系统调用把数据拷进页缓存即返回，真正的写盘由回写线程按
几十秒级的节奏异步完成。存储与设备事务相关章节讲的 DMA 与描述符环就在
回写这一刻登场：内核把脏页排进块层队列、写成描述符，磁
盘控制器按描述符用 DMA 搬数据，搬完写回完成状态，CPU
全程不经手字节。读的路径对称：命中页缓存的读是百纳秒
级，设备连被惊动都没有；未命中才排进队列、走一遍 DMA，
操作系统还会顺手预读（readahead），把顺序后继的块一并
读进页缓存。这台机器上“读文件”与“读盘”是两回事，日常
十有八九是前者。三个时延层级 $10^{2}$ ns、$10^{2}$
$\mu$s、$10$ ms 排成 $1:10^{3}:10^{5}$ 的阶梯，正是第五
章存储层次向外设方向的延伸。

设备缓存是盘内的一小块 DRAM（机械盘几十到几百 MiB，SSD
可达 GiB 级），职责有二：读预取，把顺序后继扇区提前读进
来；写暂存，把主机刚写来的数据先存下。SSD 的设备缓存还
有一份差事：“SSD 与闪存的硬件特性”一节讲的 FTL 映射表就放在里面。写暂存的
策略由 ATA 的 WCE（write cache enable）一类的开关控制。
开，对应 write-back 语义：数据一进设备缓存，控制器就向
上报完成，随后择机写介质——报完成那一刻数据还不在介质
上，断电即丢，除非盘内有掉电保护。关，对应 write-through
语义：报完成之前数据必须已在介质上，每笔写慢一个介质写
时延，换来“完成即安全”。主机侧有同名概念：页缓存本身就
是操作系统的 write-back，\code{O\_DIRECT} 与
\code{O\_SYNC} 是绕过或收紧它的两个阀门。于是“写命中哪
一层才算安全”有了明确答案：数据到达非易失的去处——介
质本身，或带掉电保护、保证能把缓存清空的盘内 RAM——这
次写才算数；停在任何一层易失 RAM 上的“完成”，都只是承
诺。

这个答案解释了企业盘与消费盘的一处规格差异。企业级 SSD
板上留一组储能电容——“电源、时钟与信号完整性”一章讲过去耦电容按时间尺度分层，
这组电容接管其中最慢的一档：市电消失——保证突然断电时
设备缓存能全部写进闪存。有了它，盘可以放心用 write-back
报完成，既快又安全；没有它的消费盘，write-back 的“完成”
在断电瞬间可能兑不了现。

两层缓存的性格也不同。页缓存面向工作集：热文件反复命中，
命中与否取决于访问的局部性；设备缓存面向流：机械盘的预
取对顺序读几乎百发百中，对随机读基本帮不上忙，于是盘
“顺序快、随机慢”的脾气在缓存层又被放大了一次。

\topic{sync/fsync 的硬件含义}
页缓存异步回写意味着 \code{write} 返回时，数据多半还躺
在主存里。要让数据真正到达介质，必须有人把每一层都走完，
这正是 \code{fsync} 的语义。一次 \code{fsync(fd)} 做三件
事：把该文件在页缓存里的脏页排进块层队列，等它们全部写
完；向设备发 FLUSH CACHE 命令，要求盘把缓存里所有暂存
写清到介质；等设备确认 FLUSH 完成，系统调用才返回。
\code{sync} 是同一件事的全盘版：清所有脏页、刷所有设备
的缓存。FLUSH 在不同协议里名字不同，角色是同一个：

\begin{longtable}{L{0.16\linewidth}L{0.34\linewidth}L{0.36\linewidth}}
\toprule
协议 & 命令名 & 语义 \\
\midrule
ATA/SATA & FLUSH CACHE（EXT） & 把盘内写缓存全部写入介质 \\\tablerowrule
SCSI/SAS & SYNCHRONIZE CACHE & 同上，可指定 LBA 区间 \\\tablerowrule
NVMe & Flush & 把指定命名空间的易失缓存写清 \\
\bottomrule
\end{longtable}

注意 FLUSH 是发给设备的命令，走设备自己的队列：在存储与设备事务相关章节
的描述符环里，它与读写命令一样占一个描述符，只是语义从
“搬数据”变成“把已搬的数据做实”。

fsync 有多贵，拆开算。极端情形先看 7200 rpm 机械盘：脏
页在盘上随机分布，每一笔都要定位——平均寻道约 4 ms，平
均旋转等待半圈 $60/7200/2\approx4.17$ ms，4 KiB 传输约
27 $\mu$s——单笔约 8-12 ms，FLUSH 再等最后一笔真正写
入磁性介质。单线程每提交一个事务就 fsync 一次，吞吐上限
被钉在 $1/10\ \mathrm{ms}=100$ 次/秒量级。换无掉电保护
的 SATA SSD：系统调用与文件系统日志处理约 10 $\mu$s，命
令下发加 DMA 搬运 4 KiB 约 10 $\mu$s，FLUSH 要等一次
TLC 页编程、约 1 ms（为什么这么慢，“SSD 与闪存的硬件特性”一节讲），合计仍是
毫秒级，上限约 $10^{3}$ 次/秒。带掉电保护的盘可以把
FLUSH 截停在缓存层——由电容提供掉电保护，缓存就算数——fsync 收
缩到几十微秒，上限抬高两个量级。一次 fsync 的时间构成
列表如下：

\begin{lstlisting}
fsync 时间构成（无掉电保护的 SATA SSD，量级）
  系统调用 + 文件系统日志写页缓存   约 10 us
  命令下发 + DMA 搬运 4 KiB        约 10 us
  FLUSH：等 TLC 页编程写入介质     约 1 ms
  合计                             毫秒级，上限约 10^3 次/秒
\end{lstlisting}

把这条层级链与 fsync 的走法画在一起，普通写与同步写
的分岔点一目了然：

\begin{pdfdiagram}
  \node[hw state, minimum width=4.6cm] (app) at (2.6,0) {应用：\code{write} / \code{fsync}};
  \node[hw memory, minimum width=4.6cm] (pc) at (2.6,-1.5) {页缓存：主存 GiB 级，读命中约 100 ns};
  \node[hw memory, minimum width=4.6cm] (dc) at (2.6,-3.0) {设备缓存：盘内 DRAM，数十 $\mu$s};
  \node[hw memory, minimum width=4.6cm] (med) at (2.6,-4.5) {介质：闪存或盘片，$10^{2}$ $\mu$s--$10$ ms};
  \draw[hw bus] (app.south) -- (pc.north);
  \draw[hw bus] (pc.south) -- (dc.north);
  \draw[hw bus] (dc.south) -- (med.north);
  \node[hw label, anchor=east] at (2.35,-0.75) {\code{write} 拷进即返回};
  \node[hw label, anchor=east] at (2.35,-2.25) {回写线程：几十秒级 DMA};
  \node[hw label, anchor=east] at (2.35,-3.75) {FLUSH 后写入介质};
  % fsync 穿透路径：右侧虚线一路向下直到介质
  \draw[hw return] (app.east) -- (6.6,0) -- (6.6,-4.5) -- (med.east);
  \node[hw label, anchor=west] at (6.8,-0.75) {① 脏页排进块层队列等写完};
  \node[hw label, anchor=west] at (6.8,-2.25) {② FLUSH 清设备缓存};
  \node[hw label, anchor=west] at (6.8,-3.75) {③ 设备确认后才返回};
  \node[hw label] at (3.3,-5.35) {普通写停在页缓存即算完成；\code{fsync} 一路穿透到介质——停在易失 RAM 的“完成”只是承诺};
\end{pdfdiagram}
\diagramnote{块设备访问的三层结构与 fsync 的穿透路径。写系统调用把数据拷进页缓存即返回，回写线程按几十秒级的节奏异步写盘，盘内还有一层设备缓存代收；fsync 则把每一层都走完：脏页排进块层队列等写完、发 FLUSH 命令清设备缓存、等设备确认才返回。数据到达非易失的去处才算数——停在任何一层易失 RAM 上的“完成”，都只是承诺。}

数据库对 fsync 延迟敏感，原因全在这笔明细里。事务日志
（WAL）的正确性规则是“日志先到介质、数据后到介质”，每
条提交记录都要过 fsync 这一关，提交时延里 fsync 常占九
成以上；它还是串行瓶颈，单个连接的提交速率就是 fsync 速
率的倒数。工程对策全在这笔明细里做文章。组提交把一批事
务攒进同一次 fsync：10 ms 内攒 200 个事务，吞吐就是
$200/10\ \mathrm{ms}=2\times10^{4}$ 事务/秒——单个事务
多等几毫秒，总量放大两百倍，“冒险、异常与性能”一章“用时延换吞吐”的交换，
这里是最贵的一例。换带掉电保护的盘，直接消除明细中代价最高
的那一项。放宽同步频率（如
\code{innodb\_flush\_log\_at\_trx\_commit=2}，每秒刷一次）
则是拿一秒的丢失窗口换两个量级。理解 fsync 的硬件含义之
后，每个调优选项都能换算成“哪一层 RAM 被信任、丢失窗口
有多大”，而不是背参数口诀。

还有一层容易漏算：文件系统自己的日志。ext4 默认
\code{data=ordered} 模式下，\code{fsync} 不只刷目标文件
的数据——若这次写伴随元数据变更（扩块、改大小），还要
先刷数据区，再等日志提交写入介质，一次 \code{fsync} 可
能触发的不止一次 FLUSH。数据库干脆自己管理日志、把数据
文件用 \code{O\_DIRECT} 打开，连页缓存这层也省掉，只留
设备缓存一层要刷——层级少一层，耗时构成就少一项不确定
性。

\section{四、从“写完成”到“掉电后仍存在”}

写系统调用返回、块请求完成、设备缓存接收、介质完成编程是四个不同边界。若应用先写数据再写指针，掉电后不能允许“新指针已持久化而新数据尚未持久化”。日志、写时复制和显式 flush 都是在安排这两个持久化顺序。

\begin{longtable}{L{0.2\linewidth}L{0.34\linewidth}L{0.36\linewidth}}
\toprule
层次 & 完成通常表示 & 尚未保证 \\
\midrule
页缓存 & 数据已被内核接收 & 块层、设备和介质完成 \\\tablerowrule
块层完成 & 设备报告请求完成 & 无断电保护缓存中的数据抗掉电 \\\tablerowrule
设备 FLUSH 完成 & 之前写入跨过设备规定的持久化边界 & 应用自己的元数据顺序必然正确 \\\tablerowrule
应用事务提交 & 日志/数据顺序满足协议 & 未纳入协议的外部系统状态 \\
\bottomrule
\end{longtable}

\section{五、校验、冗余与故障域}

ECC 修复单个器件或扇区中的有限错误，校验和发现数据路径或静默损坏，镜像与纠删码在设备失效后重建数据；它们解决的故障粒度不同。RAID 不是备份：控制器错误、软件误删、勒索软件和整机电源事故可能同时影响阵列所有成员。

分析可靠性时先画故障域：同一芯片、同一通道、同一背板、同一电源、同一机箱、同一机房。两份数据若共享同一个失效点，就不能简单按两次独立概率相乘。重建期间还会增加介质流量与暴露时间，因此容量、校验强度、备用空间和重建限速必须一起设计。
